% Section: AI Collaboration Methodology
\section{研究方法：人机协作的理论物理研究}
\label{sec:ai_methodology}

CTUFT的发展代表了理论物理研究的一种新范式，即将人类物理直觉与人工智能辅助的计算和验证相结合。本节概述本工作中采用的协作研究方法。

\subsection{分工}

研究工作流围绕"人类主导、AI辅助"的框架展开：

\begin{itemize}
    \item \textbf{人类（主导）}：负责概念框架的构建、数学结构的物理诠释、可证伪预测的设计，以及整体理论一致性把控；
    
    \item \textbf{AI（辅助）}：负责复杂的代数运算、大规模数值验证、自动化代码生成，以及全面的文献综合。
\end{itemize}

这种分工确保核心物理洞察和诠释基础始终扎根于人类的科学推理，同时利用AI的计算可扩展性和快速迭代能力。

\subsection{验证协议}

为确保可靠性和科学严谨性，所有AI生成的计算和推导都经过系统化的验证协议：

\begin{enumerate}
    \item \textbf{交叉验证}：AI生成的结果通过与独立计算方法和解析近似（如可用）进行核对；
    
    \item \textbf{物理一致性}：所有导出表达式必须满足已知的物理约束，包括量纲分析、对称性要求和守恒律；
    
    \item \textbf{极限情况分析}：通过解析可处理的极限情况测试结果，以验证渐近行为和边界条件。
\end{enumerate}

\subsection{对理论物理的启示}

本工作展示的人机协作方法，暗示了未来理论物理研究的一种变革性模式。通过将人类的创造力和物理洞察与AI处理复杂计算和模式识别的能力相结合，研究人员可以探索通过传统方法难以处理的参数空间和数学结构。这种范式并非取代人类判断，而是放大它。

\begin{insight}[人机协作研究]
AI不是理论物理学家的替代品，而是能力放大器。人类的角色从"计算执行者"转变为"物理直觉提供者"和"AI生成内容的验证者"。
\end{insight}
